\documentclass[10pt,a4paper]{amsart}
\usepackage{amsmath,amssymb,amsthm,graphicx}
% Retained minimal preamble for the Stacks Project chapter flat.tex (More on Flatness).
% Verbatim excerpt of the chapter preamble: Xy-pic setup (native preamble lines 14-18) and the native
% theorem declarations with the mathematical macros (native preamble lines 154-217). Copied mechanically;
% nothing added, removed or redefined. The wrapper class and the amsmath/amssymb/amsthm/graphicx base
% packages are supplied by the assembler, so the class/externaldocument part of the native preamble is omitted.
% For commutative diagrams we use Xy-pic
\usepackage[all]{xy}

% We use 2cell for 2-commutative diagrams.
\xyoption{2cell}
% Theorem environments.
%
\theoremstyle{plain}
\newtheorem{theorem}[subsection]{Theorem}
\newtheorem{proposition}[subsection]{Proposition}
\newtheorem{lemma}[subsection]{Lemma}

\theoremstyle{definition}
\newtheorem{definition}[subsection]{Definition}
\newtheorem{example}[subsection]{Example}
\newtheorem{exercise}[subsection]{Exercise}
\newtheorem{situation}[subsection]{Situation}

\theoremstyle{remark}
\newtheorem{remark}[subsection]{Remark}
\newtheorem{remarks}[subsection]{Remarks}

\numberwithin{equation}{subsection}

% Macros
%
\def\lim{\mathop{\mathrm{lim}}\nolimits}
\def\colim{\mathop{\mathrm{colim}}\nolimits}
\def\Spec{\mathop{\mathrm{Spec}}}
\def\Hom{\mathop{\mathrm{Hom}}\nolimits}
\def\Ext{\mathop{\mathrm{Ext}}\nolimits}
\def\SheafHom{\mathop{\mathcal{H}\!\mathit{om}}\nolimits}
\def\SheafExt{\mathop{\mathcal{E}\!\mathit{xt}}\nolimits}
\def\Sch{\mathit{Sch}}
\def\Mor{\mathop{\mathrm{Mor}}\nolimits}
\def\Ob{\mathop{\mathrm{Ob}}\nolimits}
\def\Sh{\mathop{\mathit{Sh}}\nolimits}
\def\NL{\mathop{N\!L}\nolimits}
\def\CH{\mathop{\mathrm{CH}}\nolimits}
\def\proetale{{pro\text{-}\acute{e}tale}}
\def\etale{{\acute{e}tale}}
\def\QCoh{\mathit{QCoh}}
\def\Ker{\mathop{\mathrm{Ker}}}
\def\Im{\mathop{\mathrm{Im}}}
\def\Coker{\mathop{\mathrm{Coker}}}
\def\Coim{\mathop{\mathrm{Coim}}}

% Boxtimes
%
\DeclareMathSymbol{\boxtimes}{\mathbin}{AMSa}{"02}

%
% Macros for moduli stacks/spaces
%
\def\QCohstack{\mathcal{QC}\!\mathit{oh}}
\def\Cohstack{\mathcal{C}\!\mathit{oh}}
\def\Spacesstack{\mathcal{S}\!\mathit{paces}}
\def\Quotfunctor{\mathrm{Quot}}
\def\Hilbfunctor{\mathrm{Hilb}}
\def\Curvesstack{\mathcal{C}\!\mathit{urves}}
\def\Polarizedstack{\mathcal{P}\!\mathit{olarized}}
\def\Complexesstack{\mathcal{C}\!\mathit{omplexes}}
% \Pic is the operator that assigns to X its picard group, usage \Pic(X)
% \Picardstack_{X/B} denotes the Picard stack of X over B
% \Picardfunctor_{X/B} denotes the Picard functor of X over B
\def\Pic{\mathop{\mathrm{Pic}}\nolimits}
\def\Picardstack{\mathcal{P}\!\mathit{ic}}
\def\Picardfunctor{\mathrm{Pic}}
\def\Deformationcategory{\mathcal{D}\!\mathit{ef}}

% --- appended 2026-10-12 (Chief): the project own definitions of the reference/slogan/history
% display environments, copied mechanically from the shared Stacks preamble lines 8-12
% (\usepackage{verbatim} plus the three \newenvironment lines, which the chapter preamble keeps
% verbatim from the project preamble). Nothing added, removed or redefined: these environments
% carry no mathematical content and the project itself typesets them as comments. Without these
% three lines a counted statement whose own text uses them fails with Environment slogan undefined.
\usepackage{verbatim}
\newenvironment{reference}{\comment}{\endcomment}
\newenvironment{slogan}{\comment}{\endcomment}
\newenvironment{history}{\comment}{\endcomment}
\usepackage{cleveref}
\crefname{originaldefinition}{original definition}{original definitions}
\Crefname{originaldefinition}{Original definition}{Original definitions}
\makeatletter
\newlabel{definition-big-small-h}{{0EU1}{source}{}{sourcepointer-1}{}}
\newlabel{definition-big-small-h@cref}{{[originaldefinition][0][]0EU1}{[1][0][]source}{}{}{}}
\makeatother
\crefname{originallemma}{original lemma}{original lemmas}
\Crefname{originallemma}{Original lemma}{Original lemmas}
\makeatletter
\newlabel{lemma-weak-normalization-ph-sheaf}{{0EVT}{source}{}{sourcepointer-2}{}}
\newlabel{lemma-weak-normalization-ph-sheaf@cref}{{[originallemma][0][]0EVT}{[1][0][]source}{}{}{}}
\makeatother
\crefname{originallemma}{original lemma}{original lemmas}
\Crefname{originallemma}{Original lemma}{Original lemmas}
\makeatletter
\newlabel{morphisms-lemma-seminormalization-ring}{{0EUR}{source}{}{sourcepointer-3}{}}
\newlabel{morphisms-lemma-seminormalization-ring@cref}{{[originallemma][0][]0EUR}{[1][0][]source}{}{}{}}
\makeatother
\crefname{originallemma}{original lemma}{original lemmas}
\Crefname{originallemma}{Original lemma}{Original lemmas}
\makeatletter
\newlabel{morphisms-lemma-universal-homeo-limit}{{0EUJ}{source}{}{sourcepointer-4}{}}
\newlabel{morphisms-lemma-universal-homeo-limit@cref}{{[originallemma][0][]0EUJ}{[1][0][]source}{}{}{}}
\makeatother
\title{The absolutely weakly normal structure sheaf is the sheafification of the structure sheaf for the h and ph topologies}
\author{The Stacks Project Authors}
\date{stacks-project:flat.tex:local-snapshot-20261011; GNU FDL 1.2 or later}
\begin{document}
\maketitle
\noindent\textit{Editorial note.} All delivered mathematical text is the author's own, copied from the pinned Stacks Project chapter file flat.tex; no statement or proof step was rewritten, reordered or removed.\par\medskip
\noindent\textit{Editorial note.} The preamble is the chapter's own retained minimal preamble (sources/stacks-flat/retained-preamble.tex), a verbatim excerpt of the native Xy-pic setup and of the native theorem declarations and mathematical macros, followed by the project own definitions of the reference, slogan and history display environments (verbatim package plus three newenvironment lines) that the retained file carries at its end and that the project itself typesets as comments. No macro was invented or redefined, and the wrapper class and the amsmath/amssymb/amsthm/graphicx base packages are supplied by the assembler.\par\medskip
\noindent\textit{Editorial note.} The counted statement and proof contain no cross reference, no citation and no included file, so the delivered body is self contained.\par\medskip
\noindent\textit{Editorial note.} Canonical location of the counted proposition: Stacks tag 0EVU (sources/stacks-flat/pinned-original-source.tex line 12077).\par\medskip
\noindent\textit{Allowed standard prerequisites (editorial).}\begin{itemize}
\item The absolutely weakly normal structure sheaf is a ph sheaf, with a canonical map from the structure sheaf. Same chapter, Lemma labelled lemma-weak-normalization-ph-sheaf (source line 11991), invoked at proof lines 12089-12091. Supplies that the rule is already a sheaf for ph, hence for h, and the canonical map of presheaves whose sheafification property remains to be shown.
\item The absolutely weakly normal quotient of a ring and its elements. Morphisms, Lemma labelled morphisms-lemma-seminormalization-ring (Stacks tag 0EUR), invoked at proof lines 12107-12108. Supplies that an element of F(Spec(A)) lies in A\textasciicircum{}\{awn\} and admits a finitely presented model.
\item Realising an element as the image of a finitely presented universal homeomorphism. Morphisms, Lemma labelled morphisms-lemma-universal-homeo-limit (Stacks tag 0EUJ), invoked at proof lines 12109-12113. Supplies the ring map A -> B of finite presentation with Spec(B) -> Spec(A) a universal homeomorphism over which f comes from an element of B, giving the required h covering.
\item The big and small h sites. Same chapter, Definition labelled definition-big-small-h (source line 10504), invoked in the statement. Fixes the site on which the sheafification statement is made.
\end{itemize}\medskip
\ifdefined\hypertarget\hypertarget{sourcepointer-1}{}\fi
\noindent\textit{Source reference.} Original definition 0EU1: sources/stacks-flat/pinned-original-source.tex, source lines 10504--10504. Same-chapter pointer into the author environment definition at source line 10504; invoked at proof line None of the counted proof. It is retained only as a navigation pointer; every load-bearing step is inside the delivered span.\par\medskip
\ifdefined\hypertarget\hypertarget{sourcepointer-2}{}\fi
\noindent\textit{Source reference.} Original lemma 0EVT: sources/stacks-flat/pinned-original-source.tex, source lines 11991--11991. Same-chapter pointer into the author environment lemma at source line 11991; invoked at proof line 12089 of the counted proof. It is retained only as a navigation pointer; every load-bearing step is inside the delivered span.\par\medskip
\ifdefined\hypertarget\hypertarget{sourcepointer-3}{}\fi
\noindent\textit{Source reference.} Original lemma 0EUR: sources/stacks-chow/tags.txt, source lines 18561--18561. Cross-chapter standard prerequisite located in the project tag list staged in the local snapshot; invoked at proof line 12108 of the counted proof. It is retained only as a navigation pointer; every load-bearing step is inside the delivered span.\par\medskip
\ifdefined\hypertarget\hypertarget{sourcepointer-4}{}\fi
\noindent\textit{Source reference.} Original lemma 0EUJ: sources/stacks-chow/tags.txt, source lines 18554--18554. Cross-chapter standard prerequisite located in the project tag list staged in the local snapshot; invoked at proof line 12109 of the counted proof. It is retained only as a navigation pointer; every load-bearing step is inside the delivered span.\par\medskip
\begin{lemma}
\label{lemma-weak-normalization-h-sheaf}
Let $S$ be a scheme. Choose a site $(\Sch/S)_h$
as in Definition \ref{definition-big-small-h}.
The rule
$$
X \longmapsto \Gamma(X^{awn}, \mathcal{O}_{X^{awn}})
$$
is the sheafification of the ``structure sheaf'' $\mathcal{O}$
on $(\Sch/S)_h$. Similarly for the ph topology.
\end{lemma}

\begin{proof}
In Lemma \ref{lemma-weak-normalization-ph-sheaf}
we have seen that the rule $\mathcal{F}$ of the lemma
defines a sheaf in the ph topology and hence a fortiori
a sheaf for the h topology. Clearly, there is a canonical map
of presheaves of rings $\mathcal{O} \to \mathcal{F}$.
To finish the proof, it suffices to show
\begin{enumerate}
\item if $f \in \mathcal{O}(X)$ maps to zero in $\mathcal{F}(X)$,
then there is a h covering $\{X_i \to X\}$ such that $f|_{X_i} = 0$, and
\item given $f \in \mathcal{F}(X)$ there is a h covering
$\{X_i \to X\}$ such that $f|_{X_i}$ is the image of $f_i \in \mathcal{O}(X_i)$.
\end{enumerate}
Let $f$ be as in (1). Then $f|_{X^{awn}} = 0$. This means that $f$
is locally nilpotent. Thus if $X' \subset X$ is the closed subscheme
cut out by $f$, then $X' \to X$ is a surjective closed immersion of
finite presentation. Hence $\{X' \to X\}$ is the desired h covering.
Let $f$ be as in (2). After replacing $X$ by the members of an
affine open covering we may assume $X = \Spec(A)$ is affine.
Then $f \in A^{awn}$, see
Morphisms, Lemma \ref{morphisms-lemma-seminormalization-ring}.
By Morphisms, Lemma \ref{morphisms-lemma-universal-homeo-limit}
we can find a ring map $A \to B$ of finite presentation
such that $\Spec(B) \to \Spec(A)$ is a universal homeomorphism
and such that $f$ is the image of an element $b \in B$ under
the canonical map $B \to A^{awn}$. Then
$\{\Spec(B) \to \Spec(A)\}$ is an h covering and we conclude.
The statement about the ph topology follows in the same manner
(or it can be deduced from the statement for the h topology).
\end{proof}

\end{document}
